% Clase 13 — condensadores en serie y en paralelo (primer circuitikz del curso).
% Mismos C1 y C2 en los dos paneles; cambia solo la conexion. En (a) el
% recuadro punteado es el argumento de la clase: ese tramo esta aislado y es
% neutro, asi que las dos placas que encierra llevan -Q y +Q. Adentro solo
% rotulos; lo que se comparte y lo que se suma va en el caption.
\begin{tikzpicture}[scale=1]

  % =============== (a) serie ===============
  \begin{scope}[line width=0.9pt, color=figblue]
    \draw (0,0) to[battery1] (0,2.2);
    \draw (0,2.2) to[C] (2.4,2.2) to[C] (4.8,2.2);
    \draw (4.8,2.2) -- (4.8,0) -- (0,0);
  \end{scope}
  \node[figlbl, anchor=east] at (-0.45,1.1) {$\Delta\phi$};
  \node[figlbl, figblue] at (1.2,1.3) {$C_1$};
  \node[figlbl, figblue] at (3.6,1.3) {$C_2$};
  % cargas en las placas (placas en x = 1,1 | 1,3 y 3,5 | 3,7)
  \node[figlblaux, figred, anchor=east] at (1.1,2.95) {$+Q$};
  \node[figlblaux, figblue, anchor=west] at (1.3,2.95) {$-Q$};
  \node[figlblaux, figred, anchor=east] at (3.5,2.95) {$+Q$};
  \node[figlblaux, figblue, anchor=west] at (3.7,2.95) {$-Q$};
  % tramo aislado y neutro: encierra las dos placas interiores
  \draw[figaux] (1.2,1.65) rectangle (3.6,2.75);
  \node[figlbl, anchor=north] at (2.4,-0.35) {(a) serie};

  % =============== (b) paralelo ===============
  \begin{scope}[xshift=6.6cm]
    \begin{scope}[line width=0.9pt, color=figblue]
      \draw (0,0) to[battery1] (0,2.2) -- (4.2,2.2);
      \draw (1.8,2.2) to[C] (1.8,0);
      \draw (4.2,2.2) to[C] (4.2,0);
      \draw (4.2,0) -- (0,0);
      \fill (1.8,2.2) circle (2pt);
      \fill (1.8,0) circle (2pt);
    \end{scope}
    \node[figlbl, anchor=east] at (-0.45,1.1) {$\Delta\phi$};
    \node[figlbl, figblue, anchor=east] at (1.35,1.1) {$C_1$};
    \node[figlbl, figblue, anchor=east] at (3.75,1.1) {$C_2$};
    \node[figlblaux, figred, anchor=west] at (2.2,1.45) {$+Q_1$};
    \node[figlblaux, figblue, anchor=west] at (2.2,0.75) {$-Q_1$};
    \node[figlblaux, figred, anchor=west] at (4.6,1.45) {$+Q_2$};
    \node[figlblaux, figblue, anchor=west] at (4.6,0.75) {$-Q_2$};
    \node[figlbl, anchor=north] at (2.1,-0.35) {(b) paralelo};
  \end{scope}

\end{tikzpicture}
